\chapter{常用符号表}
\label{chap:symbols}

初读可从绪论的实验开始。\secref{sec:spin-amplitudes}用自旋分析器的两个出口引入右矢、左矢与算符，
\chapref{ch:observables}再把这些记号推广到位置和动量。下表供遇到记号时回查。
帽号表示算符，粗体表示有多个空间分量，例如磁场 $\bm B$ 和自旋 $\opbm S$。

\begingroup
\small
\medskip
\renewcommand{\arraystretch}{1.22}
\noindent\begin{tabular}{@{}p{.26\linewidth}@{\hspace{1em}}p{\dimexpr.74\linewidth-1em\relax}@{}}
    \toprule
    记号 & 含义\\
    \midrule
    $\ket{\psi},\,\bra{\psi}$ & 态矢及其对偶。在正交归一基中，分别用列向量及其共轭转置表示。\\
    $\ket{a,b}$ & 逗号分隔同一个态的标签。例如 $\ket{j,m}$ 中两数分别标记角动量的大小和投影。\\
    $\ket{a}_A\otimes\ket{b}_B$ & 乘积态。约定 $a,b$ 分别表示 $A,B$ 两侧的态时，可简写为 $\ket{a,b}$。有些文献在同一约定下连写为 $\ket{ab}$（\dfnref{dfn:qm-tensor-product}）。\\
    $\braket{\phi}{\psi},\,\psi(x)$ & 内积是复数。位置波函数 $\psi(x)=\braket{x}{\psi}$ 是态在位置基中的分量。\\
    $\ketbra{a}{b},\,\op\Pi_a$ & 外积的作用是 $\ketbra{a}{b}\ket{\psi}=\ket{a}\braket{b}{\psi}$。投影 $\op\Pi_a$ 筛出读数为 $a$ 的分量，$\op\Pi_A(B)$ 筛出 $A$ 的读数落在集合 $B$ 内的分量。\\
    $\P(A=a),\,p_a$ & 正体 $\P$ 表示事件的概率。$p_a$ 是读数为 $a$ 的概率的简写。\\
    $\one,\,\one_S,\,1$ & 单位算符 $\one$ 将每个态矢映回自身。$\one_S\otimes\op A_D$ 只对仪器 $D$ 作变换。$1$ 表示数，$I$ 可表示电流。\\
    $\op A\adjoint,\,[\op A,\op B]$ & 前者是伴随算符。对易子 $[\op A,\op B]=\op A\op B-\op B\op A$ 比较两种作用次序。$\op A\op B\ket{\psi}$ 中先作用 $\op B$。\\
    $\op X,\,\op P,\,x,\,p$ & 一维位置、动量算符和相应的读数。$\opbm X$ 表示由三个位置分量算符组成的矢量。\\
    $\ave{\op A},\,\Delta A$ & 同一制备下，测量 $A$ 的读数均值和标准差。$\Delta A$ 与读数单位相同（\dfnref{dfn:qm-measurement-moments}）。\\
    $\op\rho,\,\rho_{eg},\,\mathcal E$ & $\op\rho$ 是密度算符，$\rho_{eg}=\mel{e}{\op\rho}{g}$ 是其矩阵元。量子信道 $\mathcal E$ 把输入的密度算符映为输出的密度算符。\\
    $\op U(t,t_0),\,\hbar,\,\omega$ & $\op U(t,t_0)$ 把 $t_0$ 的态演化到 $t$。$\hbar=h/(2\pi)$，角频率 $\omega=2\pi\nu$，光子能量为 $\hbar\omega=h\nu$。\\
    \bottomrule
\end{tabular}\par

\medskip
\noindent\textbf{各章中 $A$ 的几种常见含义}

\smallskip
\noindent\begin{tabular}{@{}p{.26\linewidth}@{\hspace{1em}}p{\dimexpr.74\linewidth-1em\relax}@{}}
    \toprule
    记号与出处 & 物理含义\\
    \midrule
    $\opbm A$，\chapref{ch:hydrogen} & Runge--Lenz 矢量算符。它的三个分量都作用于轨道态。\\
    $A_{e\to g}$，\chapref{ch:transitions} & 从 $e$ 到 $g$ 的自发发射率，单位为 $\unit{\per\second}$。\\
    $\bm A(\bm x,t)$，\chapref{ch:em} & 经典磁矢势的场值。$\bm A(\opbm X,t)$ 是相应的乘法算符，\chapref{ch:relativity}将它记为 $\opbm A$。\\
    $A(a,\lambda)$，\chapref{ch:entanglement} & 在 Bell 实验的局域模型中，给定设置 $a$ 和隐藏变量 $\lambda$ 后，左侧仪器的输出为 $+1$ 或 $-1$。\\
    $\bm{\mathcal A}_n$，\chapref{ch:adiabatic} & 第 $n$ 个本征态的 Berry 连接。沿外参量闭路积分给出几何相位（模 $2\pi$，沿回路取单值规范）。\\
    $A(t)$，\chapref{ch:open} & 生存振幅，即演化后的态与初态的重叠。其模平方是生存概率。\\
    $A$，\chapref{ch:relativity} & 超精细相互作用的耦合常量，具有能量量纲。此处 $\opbm I$ 是核自旋矢量算符，$I$ 是核自旋量子数。\\
    \bottomrule
\end{tabular}\par
\endgroup
